% \subsection{Обозначения (читать первым)}
% \textbf{Факториальные степени.} $x^{\underline k}=x(x-1)(x-2)\cdots(x-k+1)$ -- \emph{убывающая} степень ($k$ множителей, каждый на 1 меньше предыдущего); $x^{\underline k}=k!\binom xk$, это число упорядоченных выборок $k$ элементов из $x$. $x^{\overline k}=x(x+1)\cdots(x+k-1)$ -- \emph{возрастающая}. Чёрточка под показателем -- убывающая, над ним -- возрастающая.
% \par $[P]$ -- скобка Айверсона: $1$, если утверждение $P$ верно, иначе $0$ (например $[n=1]$). $\lfloor x\rfloor$, $\lceil x\rceil$ -- пол и потолок, $\{x\}=x-\lfloor x\rfloor$ -- дробная часть. $[x^k]F(x)$ -- коэффициент при $x^k$ в ряде $F$. $a\bmod m$ -- остаток в $[0,m)$.
% % \par \textbf{Числа.} $\binom nk$ (равно 0 вне $0\le k\le n$); $H_n=1+\frac12+\dots+\frac1n$; $F_n$ -- Фибоначчи ($F_0{=}0,F_1{=}1$); $C_n$ -- Каталан; $D_n$ -- беспорядки (перестановки без неподвижных точек); $S(n,k)$ и $c(n,k)$ -- Стирлинга II и I рода; $B_n$ -- Белла; $p(n)$ -- число разбиений.
% \par \textbf{Теория чисел.} $\varphi$ -- Эйлер, $\mu$ -- Мёбиус, $d(n)$ (иногда $\tau$) -- число делителей, $\sigma(n)$ -- их сумма, $v_p(n)$ -- показатель простого $p$ в $n$, $a\perp b$ значит $\gcd(a,b)=1$. $f*g$ -- свёртка Дирихле: $(f*g)(n)=\sum_{d\mid n}f(d)g(n/d)$; функции $\mathbf1(n)=1$, $\mathrm{id}(n)=n$, $\varepsilon(n)=[n=1]$.
% \par \textbf{Векторы.} $u\cdot v=u_xv_x+u_yv_y$ -- скалярное, $u\times v=u_xv_y-u_yv_x$ -- векторное (псевдоскалярное): знак $>0$, если $v$ повёрнут против часовой от $u$; $|u\times v|$ -- площадь параллелограмма.
% \par \textbf{Графы.} $\alpha$ -- макс. независимое множество, $\nu$ -- макс. паросочетание, $\beta$ -- мин. вершинное покрытие, $\rho$ -- мин. рёберное покрытие, $\chi$ -- хроматическое число, $\mathrm{odd}(G)$ -- число компонент нечётного размера. $\oplus$ -- xor, $\mathrm{mex}(A)$ -- наименьшее неотрицательное число не из $A$. $E[X]$, $V(X)$ -- матожидание и дисперсия.
% \par \textbf{Производящие функции.} Обычная (OGF): $\sum a_nx^n$. Экспоненциальная (EGF): $\sum a_n\frac{x^n}{n!}$. EGF удобна для \emph{помеченных} объектов (перестановки, разбиения), OGF -- для \emph{непомеченных} (композиции, разбиения чисел).
